% ============================================================
%  SCRIPT NGƯỜI 2 — INGESTION + REDPANDA + PREFECT
%  Olist Data Lakehouse Platform — 3 phút (3:00 -- 6:00)
%  Compile: pdflatex -interaction=nonstopmode script_nguoi2.tex
% ============================================================
\documentclass[12pt, a4paper]{article}

\usepackage[utf8]{inputenc}
\usepackage[T5]{fontenc}
\usepackage[vietnamese]{babel}
\usepackage[top=2.2cm, bottom=2.2cm, left=2.5cm, right=2.5cm]{geometry}
\usepackage{parskip}
\setlength{\parindent}{0pt}
\usepackage{lmodern}
\usepackage{microtype}
\usepackage{setspace}
\onehalfspacing

% ── Colors ───────────────────────────────────────────────────
\usepackage[dvipsnames, table, x11names]{xcolor}
\definecolor{p2color}  {RGB}{5, 150, 105}
\definecolor{bronze}   {RGB}{180, 103, 40}
\definecolor{silver}   {RGB}{130, 140, 150}
\definecolor{gold}     {RGB}{180, 145, 20}
\definecolor{demoboxbg}{RGB}{15,  23,  42}
\definecolor{scriptbg} {RGB}{241, 245, 249}
\definecolor{titlebg}  {RGB}{15,  23,  42}
\definecolor{warnbg}   {RGB}{255, 251, 235}
\definecolor{codebg}   {RGB}{30,  41,  59}

% ── Boxes ────────────────────────────────────────────────────
\usepackage[most]{tcolorbox}
\tcbuselibrary{breakable, skins, listings}

% Speaker box
\newenvironment{speakerbox}[2]{%
  \begin{tcolorbox}[
    enhanced, breakable,
    colback=#1!6!white, colframe=#1,
    colbacktitle=#1, coltitle=white,
    fonttitle=\bfseries\large,
    title={#2},
    arc=4pt, boxrule=1.5pt,
    left=8pt, right=8pt, top=6pt, bottom=6pt,
  ]
}{\end{tcolorbox}}

% Demo action box
\newtcolorbox{demobox}[1][]{
  enhanced, breakable,
  colback=demoboxbg, colframe=cyan!60,
  coltext=white,
  fonttitle=\bfseries\small\color{cyan!80},
  title={DEMO --- Thao tác màn hình},
  arc=4pt, boxrule=1pt,
  left=8pt, right=8pt, #1
}

% Script / speech style
\newtcolorbox{scriptbox}{
  enhanced, breakable,
  colback=scriptbg, colframe=gray!40,
  arc=3pt, boxrule=0.8pt,
  left=12pt, right=8pt, leftrule=4pt,
}

% Code box
\newtcolorbox{codebox}[1]{
  enhanced, breakable,
  colback=codebg, colframe=p2color!60,
  coltext=white, coltitle=white,
  fonttitle=\bfseries\small,
  title={#1},
  arc=3pt, boxrule=1pt,
  left=8pt, right=8pt,
}

% Key point box
\newtcolorbox{keybox}[1]{
  enhanced,
  colback=p2color!8, colframe=p2color!50,
  coltitle=p2color, fonttitle=\bfseries\small,
  title={#1},
  arc=3pt, boxrule=1pt,
  left=8pt, right=8pt,
}

% Warning / prep box
\newtcolorbox{warnbox}[1]{
  enhanced,
  colback=warnbg, colframe=orange!60,
  coltitle=orange!70!black, fonttitle=\bfseries\small,
  title={#1},
  arc=3pt, boxrule=1pt,
  left=8pt, right=8pt,
}

% ── Helpers ──────────────────────────────────────────────────
\newcommand{\timing}[1]{\hfill{\small\bfseries\color{gray}[#1]}}
\newcommand{\hl}[1]{\textbf{\color{p2color}#1}}
\newcommand{\code}[1]{\texttt{\small#1}}

% ── TikZ ─────────────────────────────────────────────────────
\usepackage{tikz}
\usetikzlibrary{arrows.meta, positioning, calc, shapes.geometric, fit, backgrounds}
\usepackage{fontawesome5}
\usepackage{booktabs}
\usepackage{tabularx}
\usepackage{enumitem}
\usepackage{multicol}

% ── Hyperref ─────────────────────────────────────────────────
\usepackage{hyperref}
\hypersetup{colorlinks=true, linkcolor=p2color, urlcolor=p2color}

% ── Header / Footer ──────────────────────────────────────────
\usepackage{fancyhdr}
\pagestyle{fancy}
\fancyhf{}
\fancyhead[L]{\small\color{gray}Script Người 2 --- Ingestion + Redpanda + Prefect}
\fancyhead[R]{\small\color{gray}3:00 -- 6:00 \textbullet\ 3 phút}
\fancyfoot[C]{\small\color{gray}\thepage}
\renewcommand{\headrulewidth}{0.4pt}

% ============================================================
\begin{document}
% ============================================================

% ── Title ────────────────────────────────────────────────────
\begin{tikzpicture}[remember picture, overlay]
  \fill[titlebg]
    (current page.north west) rectangle
    ([yshift=-5.5cm]current page.north east);
  \node[anchor=north, text=white, font=\LARGE\bfseries, yshift=-1.5cm]
    at (current page.north)
    {Olist Data Lakehouse --- Script Người 2};
  \node[anchor=north, text=p2color!80, font=\large, yshift=-2.8cm]
    at (current page.north)
    {Ingestion Layer \textbullet\ Redpanda Streaming \textbullet\ Prefect Orchestration};
  \node[anchor=north, text=white!60, font=\normalsize, yshift=-4.0cm]
    at (current page.north)
    {Thời gian: 3:00 -- 6:00 \textbullet\ 3 phút \textbullet\ 2 demo live};
\end{tikzpicture}

\vspace{6cm}

\begin{warnbox}{\faExclamationTriangle\ Chuẩn bị trước khi lên trình bày}
\begin{itemize}[leftmargin=1.5em, itemsep=3pt]
  \item \textbf{Tab Redpanda Console} mở sẵn tại \code{http://localhost:8080}
        --- topic \code{olist.orders} đã có message
  \item \textbf{Tab Prefect UI} mở sẵn tại \code{http://localhost:4200}
        --- có ít nhất 1 flow run \textcolor{green!60!black}{Completed}
  \item Producer đã chạy: \code{docker exec olist-producer python streaming/producer.py}
  \item Consumer đã chạy: \code{docker exec olist-consumer python streaming/consumer.py}
  \item Biết sẵn câu trả lời: ``Tại sao dùng Redpanda mà không dùng Kafka thật?''
\end{itemize}
\end{warnbox}

\vspace{0.5cm}

\begin{center}
\begin{tikzpicture}[x=1cm, y=0.55cm]
  \draw[line width=2.5pt, gray!30] (0,0) -- (15,0);
  \fill[p2color!80, rounded corners=2pt] (3,-0.28) rectangle (6,0.28);
  \fill[gray!30,   rounded corners=2pt] (0,-0.28) rectangle (2.95,0.28);
  \fill[gray!30,   rounded corners=2pt] (6.05,-0.28) rectangle (15,0.28);
  \node[text=white, font=\small\bfseries] at (4.5, 0) {Người 2 --- 3'};
  \node[font=\tiny, gray] at (1.5, 0)  {Người 1};
  \node[font=\tiny, gray] at (7.5, 0)  {Người 3};
  \node[font=\tiny, gray] at (12,  0)  {Người 4};
  \foreach \x in {0,3,...,15} {
    \draw[gray!50] (\x, -0.5) -- (\x, -0.7);
    \node[font=\tiny, gray] at (\x, -1.0) {\x'};
  }
\end{tikzpicture}
\end{center}

\newpage
\tableofcontents
\newpage

% ============================================================
\section{Phân bố thời gian chi tiết (3 phút)}
% ============================================================

\begin{center}
\renewcommand{\arraystretch}{1.6}
\begin{tabularx}{\textwidth}{clXl}
\toprule
\textbf{Giai đoạn} & \textbf{Thời điểm} & \textbf{Nội dung} & \textbf{Loại} \\
\midrule
0 & 3:00 -- 3:10 & Tiếp nhận mic từ Người 1, giới thiệu nhanh phần mình & Lời nói \\
1 & 3:10 -- 3:40 & Batch ingestion --- CSV $\to$ PyIceberg Bronze & Lời nói + code \\
2 & 3:40 -- 4:30 & Streaming --- Redpanda + producer/consumer (DEMO) & Demo live \\
3 & 4:30 -- 5:20 & Prefect UI --- flow run + task graph + quality gate (DEMO) & Demo live \\
4 & 5:20 -- 5:50 & Tổng kết Ingestion layer + chuyển tiếp & Lời nói \\
5 & 5:50 -- 6:00 & \textit{Bước đệm, chuẩn bị Người 3 lên} & \textit{Transition} \\
\bottomrule
\end{tabularx}
\end{center}

\begin{keybox}{\faLightbulb\ Mục tiêu của phần này}
Người xem hiểu \textbf{hai luồng vào cùng một Bronze layer}: batch qua PyIceberg và
streaming qua Redpanda, không conflict nhờ ACID của Iceberg.
Và thấy \textbf{Prefect tự động hóa} toàn bộ pipeline với quality gate.
\end{keybox}

\newpage

% ============================================================
\section{Giai đoạn 0 --- Tiếp nhận \& Mở đầu (3:00 -- 3:10)}
% ============================================================

\begin{speakerbox}{p2color}{\faUser\ Người 2 \textbullet\ 3:00 -- 3:10 \textbullet\ Tiếp nhận}

\begin{scriptbox}
``Cảm ơn bạn. Vậy là chúng ta đã hiểu cách system nhìn tổng thể. Bây giờ phần của
mình sẽ đi sâu vào \textbf{hai đầu vào của hệ thống} --- dữ liệu vào Bronze
layer bằng cách nào, và ai điều phối toàn bộ flow.''
\end{scriptbox}

\end{speakerbox}

% ============================================================
\section{Giai đoạn 1 --- Batch Ingestion (3:10 -- 3:40)}
% ============================================================

\begin{speakerbox}{p2color}{\faUser\ Người 2 \textbullet\ 3:10 -- 3:40 \textbullet\ Lời nói}

\subsection*{Nội dung cần nói}

\begin{scriptbox}
``Luồng đầu tiên --- \textbf{Batch Ingestion}. 8 file CSV gốc của Olist,
tổng khoảng 200~MB, được đọc bằng Pandas và ghi vào
\textbf{Bronze layer trên Cloudflare R2} thông qua PyIceberg.

Thiết kế quan trọng nhất của Bronze: \textbf{append-only, immutable}.
Tức là không bao giờ UPDATE hay DELETE. Nếu source data thay đổi,
chúng ta \textit{append} snapshot mới. Lịch sử luôn được giữ
nguyên --- đây là nền tảng của time travel sau này.

Kết quả: 10 Iceberg table trên R2, mỗi table có \texttt{data/} chứa
Parquet và \texttt{metadata/} chứa snapshots, manifests --- đây là cấu trúc
của Iceberg Table Format.''
\end{scriptbox}

\end{speakerbox}

\vspace{0.5em}

\begin{codebox}{\faCode\ prefect/flows/bronze\_ingestion.py --- đoạn key}
\begin{verbatim}
@task(retries=3, retry_delay_seconds=30)
def ingest_dataset(dataset: str, catalog) -> int:
    table = catalog.load_table(f"bronze.ecom.{dataset}")
    df = pd.read_csv(f"data/olist_{dataset}_dataset.csv")
    df = df.where(pd.notna(df), None)   # NaN -> None (Iceberg-safe)

    schema = table.schema()
    pa_table = pa.Table.from_pandas(df, schema=schema.as_arrow())
    table.append(pa_table)   # <-- append-only, ACID
    return len(df)

@flow(name="bronze-batch-ingestion")
def bronze_batch_flow():
    catalog = load_catalog("olist", **CATALOG_CONF)
    datasets = ["orders", "order_items", "customers",
                "sellers", "products", "payments",
                "reviews", "geolocation"]
    for ds in datasets:
        count = ingest_dataset(ds, catalog)
        logger.info(f"Ingested {count} rows -> bronze.ecom.{ds}")
\end{verbatim}
\end{codebox}

\begin{keybox}{\faKey\ Điểm kỹ thuật cần nhấn mạnh}
\begin{itemize}[leftmargin=1.5em, itemsep=4pt]
  \item \code{@task(retries=3)} --- nếu S3/R2 timeout, tự retry, không cần can thiệp thủ công
  \item \code{table.append(pa\_table)} --- PyIceberg API, không bao giờ dùng raw Parquet write
  \item \code{df.where(pd.notna(df), None)} --- chuyển NaN sang None, Iceberg không hiểu NaN của pandas
  \item Bronze table schema đã được tạo sẵn qua PyIceberg schema definition
\end{itemize}
\end{keybox}

\newpage

% ============================================================
\section{Giai đoạn 2 --- Redpanda Streaming (3:40 -- 4:30)}
% ============================================================

\begin{speakerbox}{p2color}{\faUser\ Người 2 \textbullet\ 3:40 -- 4:30 \textbullet\ DEMO + Lời nói}

\subsection*{[Bước 1] Mở Redpanda Console \timing{3:40 -- 3:55}}

\begin{demobox}
\begin{enumerate}[leftmargin=1.5em, itemsep=4pt]
  \item Switch sang tab \textbf{Redpanda Console} (localhost:8080)
  \item Click menu bên trái: \textbf{Topics} $\to$ click vào \code{olist.orders}
  \item Thấy danh sách messages --- click một message bất kỳ để xem payload JSON
  \item Expand JSON ra để thấy \code{order\_id}, \code{customer\_id}, \code{timestamp}
\end{enumerate}
\end{demobox}

\begin{scriptbox}
``Đây là \textbf{Redpanda Console} --- giao diện giám sát broker. Topic
\code{olist.orders} đang chạy real-time. Nhìn vào một message bất kỳ ---
\textbf{payload JSON} có \code{order\_id}, \code{customer\_id}, \code{timestamp}.

Message key là \texttt{order\_id} --- điều này đảm bảo tất cả sự kiện
của cùng một đơn hàng rơi vào cùng một partition --- \textbf{ordering
guarantee} cho mỗi đơn hàng.''
\end{scriptbox}

\subsection*{[Bước 2] Giải thích kiến trúc Streaming \timing{3:55 -- 4:15}}

\begin{scriptbox}
``Redpanda chạy trong một container duy nhất --- không cần Zookeeper, không
cần JVM. RAM tiêu thụ dưới \textbf{500~MB}, so với Kafka thật cần \textbf{2~GB}.
API hoàn toàn Kafka-compatible --- code Python dùng \code{confluent-kafka} không cần
thay đổi một dòng.

\textbf{Producer} (\code{streaming/producer.py}) đọc CSV gốc, sort theo
\code{order\_purchase\_timestamp}, rồi replay theo thời gian thật với
\textbf{speed factor 86400} --- tức là 1 ngày lịch sử bằng 1 giây
real-time. Điều này giúp demo streaming mà không cần chờ thật.

\textbf{Consumer} (\code{streaming/consumer.py}) thuộc consumer group
\code{iceberg-bronze-writer} --- nhận message và ghi vào cùng Bronze Iceberg
table với \code{table.append()} --- cùng API với batch, không conflict nhờ
ACID của Iceberg.''
\end{scriptbox}

\subsection*{[Bước 3] Giải thích Consumer code \timing{4:15 -- 4:30}}

\begin{scriptbox}
``Consumer chạy liên tục. Mỗi lần nhận được 50 message, nó flush
một lần vào Iceberg --- mô hình \textbf{micro-batch} giúp giảm overhead ghi
và tăng throughput. Nếu consumer restart giữa chừng, Kafka offset đã
commit --- không mất message, không duplicate.''
\end{scriptbox}

\end{speakerbox}

\vspace{0.5em}

\begin{codebox}{\faCode\ streaming/producer.py --- Speed replay logic}
\begin{verbatim}
SPEED_FACTOR = 86_400   # 1 day of history = 1 second realtime

df = pd.read_csv("data/olist_orders_dataset.csv",
                 parse_dates=["order_purchase_timestamp"])
events = df.sort_values("order_purchase_timestamp").itertuples()

prev_ts = None
for ev in events:
    if prev_ts is not None:
        gap = (ev.order_purchase_timestamp - prev_ts
               ).total_seconds() / SPEED_FACTOR
        if gap > 0:
            time.sleep(gap)
    producer.produce(
        topic="olist.orders",
        key=ev.order_id,                     # partition ordering
        value=json.dumps({
            "event_type": "order_created",
            "order_id":   ev.order_id,
            "customer_id": ev.customer_id,
            "timestamp":  str(ev.order_purchase_timestamp),
        }),
    )
    producer.poll(0)
    prev_ts = ev.order_purchase_timestamp
producer.flush()
\end{verbatim}
\end{codebox}

\begin{codebox}{\faCode\ streaming/consumer.py --- Micro-batch flush to Iceberg}
\begin{verbatim}
BATCH_SIZE = 50
consumer = Consumer({
    "bootstrap.servers": "redpanda:9092",
    "group.id": "iceberg-bronze-writer",    # consumer group
    "auto.offset.reset": "earliest",
})
consumer.subscribe(["olist.orders"])

buffer = []
while True:
    msg = consumer.poll(timeout=1.0)
    if msg is None:
        continue
    payload = json.loads(msg.value())
    buffer.append(payload)

    if len(buffer) >= BATCH_SIZE:
        pa_batch = pa.Table.from_pylist(buffer, schema=BRONZE_SCHEMA)
        iceberg_table.append(pa_batch)   # ACID append to Bronze
        consumer.commit()                 # commit offset AFTER flush
        buffer.clear()
\end{verbatim}
\end{codebox}

\vspace{0.5em}

% ── Diagram: Streaming Architecture ──────────────────────────
\begin{tcolorbox}[colback=gray!5, colframe=p2color!50, arc=6pt,
  title={\faSitemap\ Sơ đồ --- Luồng Streaming vào Bronze Iceberg}]
\begin{center}
\begin{tikzpicture}[
  node distance=0.5cm and 0.9cm,
  box/.style={
    draw, rounded corners=4pt, minimum width=2.4cm,
    minimum height=0.75cm, align=center, font=\small, line width=1pt
  },
  arr/.style={-Stealth, line width=1.2pt},
]

\node[box, fill=gray!15, draw=gray] (csv)
  {CSV Orders\\{\tiny 100K rows, sorted by ts}};

\node[box, fill=orange!20, draw=orange!70, right=1.0cm of csv] (prod)
  {\textbf{Producer}\\{\tiny speed factor 86400}};

\node[box, fill=orange!35, draw=orange, right=1.0cm of prod,
      minimum height=1.3cm] (rp)
  {\faRss\ \textbf{Redpanda}\\{\tiny topic: olist.orders}\\{\tiny key = order\_id}};

\node[box, fill=green!20, draw=green!60, right=1.0cm of rp,
      minimum height=1.3cm] (cons)
  {\textbf{Consumer}\\{\tiny group: iceberg-bronze-writer}\\{\tiny micro-batch=50}};

\node[box, fill=orange!25, draw=orange!80, right=1.0cm of cons,
      minimum height=1.3cm] (bronze)
  {\textbf{BRONZE}\\{\tiny Iceberg table}\\{\tiny R2 / S3}};

\draw[arr, gray]       (csv)  -- (prod);
\draw[arr, orange!80]  (prod) -- node[above, font=\tiny]{produce} (rp);
\draw[arr, green!60]   (rp)   -- node[above, font=\tiny]{poll} (cons);
\draw[arr, orange!80]  (cons) -- node[above, font=\tiny]{append} (bronze);

\node[font=\tiny, color=gray,      below=0.1cm of cons]   {commit offset after flush};
\node[font=\tiny, color=orange!80, below=0.1cm of bronze]  {ACID \textbullet\ append-only};
\node[font=\tiny, color=gray,      above=0.1cm of prod]    {1 day history = 1s realtime};

\end{tikzpicture}
\end{center}

\begin{center}
\small
\begin{tabular}{lll}
\toprule
\textbf{Topic} & \textbf{Key} & \textbf{Consumer Group} \\
\midrule
\code{olist.orders}  & \code{order\_id}  & \code{iceberg-bronze-writer} \\
\code{olist.reviews} & \code{review\_id} & \code{iceberg-bronze-writer} \\
\code{olist.leads}   & \code{mql\_id}    & \code{iceberg-bronze-writer} \\
\code{olist.deals}   & \code{mql\_id}    & \code{iceberg-bronze-writer} \\
\bottomrule
\end{tabular}
\end{center}
\end{tcolorbox}

\newpage

% ============================================================
\section{Giai đoạn 3 --- Prefect Orchestration (4:30 -- 5:20)}
% ============================================================

\begin{speakerbox}{p2color}{\faUser\ Người 2 \textbullet\ 4:30 -- 5:20 \textbullet\ DEMO Prefect UI}

\subsection*{[Bước 1] Vào Prefect UI, chọn flow run \timing{4:30 -- 4:45}}

\begin{demobox}
\begin{enumerate}[leftmargin=1.5em, itemsep=4pt]
  \item Switch sang tab \textbf{Prefect UI} (localhost:4200)
  \item Click \textbf{Flow Runs} trên sidebar trái
  \item Tìm một run tên \code{full-pipeline} có status \textcolor{green!60!black}{\faCheckCircle\ Completed}
  \item Click vào run đó để mở detail view
\end{enumerate}
\end{demobox}

\begin{scriptbox}
``Đây là \textbf{Prefect UI} --- nơi giám sát toàn bộ pipeline. Mỗi
hình tròn trên \textbf{Task Graph} này là một task Python. Màu xanh lá là
\textcolor{green!60!black}{Completed}, màu đỏ là Failed.

Thấy rõ thứ tự: \textbf{Bronze Ingestion} $\to$ \textbf{Soda Check Bronze}
$\to$ \textbf{Silver Transform} $\to$ \textbf{Soda Check Silver} $\to$
\textbf{Gold Transform} --- toàn bộ luồng được định nghĩa bằng
\code{wait\_for} trong Python, không cần YAML, không cần công cụ riêng.''
\end{scriptbox}

\subsection*{[Bước 2] Click vào Soda Check task \timing{4:45 -- 5:05}}

\begin{demobox}
\begin{enumerate}[leftmargin=1.5em, itemsep=4pt]
  \item Trên Task Graph, click vào task có tên \code{soda\_check\_bronze}
  \item Click tab \textbf{Logs} để xem output của Soda Core
  \item Đọc một vài dòng log: \textit{``Pass: not\_null check on order\_id...''} v.v.
  \item Quay lại Task Graph, chỉ vào mũi tên giữa Soda và Silver
\end{enumerate}
\end{demobox}

\begin{scriptbox}
``Task \textbf{Soda Check} là \textbf{quality gate} giữa mỗi layer. Soda Core kiểm tra:

Thứ nhất là \textbf{null rate} --- nếu \code{order\_id} có null, fail ngay.
Thứ hai là \textbf{row count} --- phải có ít nhất 90 nghìn dòng.
Thứ ba là \textbf{referential integrity} --- mỗi order phải tồn tại trong bảng
customers.

Nếu bất kỳ check nào fail, Prefect \textbf{dừng flow lại ngay tại
đây}, không tiếp tục lên Silver. Và gửi webhook notification ---
thực tế có thể dùng Slack hay email. Đây là \textbf{data quality as
code}, không cần kiểm tra thủ công.''
\end{scriptbox}

\subsection*{[Bước 3] Chỉ Deployments \timing{5:05 -- 5:20}}

\begin{demobox}
\begin{enumerate}[leftmargin=1.5em, itemsep=3pt]
  \item Click \textbf{Deployments} trên sidebar
  \item Chỉ vào schedule \code{Daily 02:00 UTC} của \code{full-pipeline}
  \item Chỉ nhanh 4 flow có trong hệ thống
\end{enumerate}
\end{demobox}

\begin{scriptbox}
``Trong \textbf{Deployments}, mỗi flow đã được đăng ký với schedule.
\code{full-pipeline} chạy 2 giờ sáng mỗi ngày. \code{ml-training} chạy hàng
tuần để retrain model. \code{sentiment-flow} chạy 3 giờ sáng score các
review mới.

Quan trọng: Prefect Worker chạy trong container
\code{prefect-worker}, có đủ Python + Spark. Nếu worker bị restart,
các scheduled flow vẫn chạy bình thường --- Prefect Server giữ
state riêng.''
\end{scriptbox}

\end{speakerbox}

\vspace{0.5em}

\begin{codebox}{\faCode\ prefect/flows/full\_pipeline.py --- Flow definition}
\begin{verbatim}
from prefect import flow, task
from prefect.logging import get_run_logger

@task(retries=3, retry_delay_seconds=60)
def bronze_ingestion_task():
    from bronze_ingestion import bronze_batch_flow
    bronze_batch_flow()

@task(retries=2)
def soda_check_task(layer: str) -> bool:
    from quality.soda_runner import run_checks
    result = run_checks(layer)          # returns True if all pass
    if not result:
        raise ValueError(f"Soda check FAILED for layer={layer}")
    return True

@task(retries=2, retry_delay_seconds=30)
def spark_transform_task(job: str):
    import subprocess
    subprocess.run(
        ["spark-submit", f"spark/jobs/{job}.py"],
        check=True                      # raise on non-zero exit
    )

@flow(name="full-pipeline", log_prints=True)
def full_pipeline():
    logger = get_run_logger()

    logger.info("Step 1: Bronze ingestion")
    bronze_ingestion_task()

    logger.info("Step 2: Soda check Bronze")
    soda_check_task("bronze")           # blocks if fail

    logger.info("Step 3: Silver transform")
    spark_transform_task("silver_transform")

    logger.info("Step 4: Soda check Silver")
    soda_check_task("silver")

    logger.info("Step 5: Gold transform")
    spark_transform_task("gold_transform")

    logger.info("Pipeline complete!")
\end{verbatim}
\end{codebox}

\vspace{0.5em}

% ── Diagram: Prefect Flow ─────────────────────────────────────
\begin{tcolorbox}[colback=gray!5, colframe=p2color!50, arc=6pt,
  title={\faSitemap\ Sơ đồ --- Prefect Full Pipeline với Quality Gates}]
\begin{center}
\begin{tikzpicture}[
  node distance=0.4cm and 0.8cm,
  task/.style={
    draw, rounded corners=4pt, minimum width=2.3cm,
    minimum height=0.65cm, align=center, font=\small, line width=1pt
  },
  gate/.style={
    diamond, draw, aspect=2.0, minimum width=2.6cm,
    minimum height=0.5cm, align=center, font=\footnotesize,
    fill=yellow!20, draw=orange!60, line width=1pt
  },
  fail/.style={
    task, fill=red!15, draw=red!50, font=\footnotesize
  },
  arr/.style={-Stealth, line width=1pt},
  lbl/.style={font=\tiny, color=gray},
]

\node[task, fill=gray!15, draw=gray]         (start)  {Start};
\node[task, fill=orange!25, draw=orange!70,
      right=of start]                         (bronze) {Bronze\\Ingestion};
\node[gate, right=of bronze]                 (qb)     {Soda\\Bronze?};
\node[fail, below=0.65cm of qb]             (failb)  {PAUSE\\+ notify};
\node[task, fill=silver!40, draw=silver!80,
      right=of qb]                            (silver) {Silver\\Transform};
\node[gate, right=of silver]                (qs)     {Soda\\Silver?};
\node[fail, below=0.65cm of qs]             (fails)  {PAUSE\\+ notify};
\node[task, fill=gold!40, draw=gold!90,
      right=of qs]                            (gold)   {Gold\\Transform};
\node[task, fill=green!15, draw=green!60,
      right=of gold]                          (done)   {\faCheck\ Done};

\draw[arr]              (start)  -- (bronze);
\draw[arr]              (bronze) -- (qb);
\draw[arr]              (qb)     -- node[lbl,above]{Pass} (silver);
\draw[arr, red!50]      (qb)     -- node[lbl,right]{Fail} (failb);
\draw[arr]              (silver) -- (qs);
\draw[arr]              (qs)     -- node[lbl,above]{Pass} (gold);
\draw[arr, red!50]      (qs)     -- node[lbl,right]{Fail} (fails);
\draw[arr]              (gold)   -- (done);

\draw[dashed, orange!60, -Stealth, thin]
  (failb.west) to[bend right=50]
  node[lbl,below]{\tiny retry 3x} (bronze.south);

\node[lbl, above=0.1cm of bronze] {\tiny retries=3};
\node[lbl, above=0.1cm of silver] {\tiny retries=2};
\node[lbl, above=0.1cm of gold]   {\tiny retries=2};

\end{tikzpicture}
\end{center}

\begin{center}
\small
\begin{tabular}{lll}
\toprule
\textbf{Flow} & \textbf{Mô tả} & \textbf{Schedule} \\
\midrule
\code{full-pipeline}  & Bronze $\to$ Silver $\to$ Gold + quality gates & Daily 02:00 UTC \\
\code{ml-training}    & Retrain 3 XGBoost models                        & Weekly \\
\code{sentiment-flow} & Score reviews mới bằng BERTimbau                & Daily 03:00 UTC \\
\code{quality-checks} & Standalone Soda check                            & On-demand \\
\bottomrule
\end{tabular}
\end{center}
\end{tcolorbox}

\newpage

% ============================================================
\section{Giai đoạn 4 --- Tổng kết \& Chuyển tiếp (5:20 -- 6:00)}
% ============================================================

\begin{speakerbox}{p2color}{\faUser\ Người 2 \textbullet\ 5:20 -- 6:00 \textbullet\ Tổng kết}

\subsection*{Tóm tắt Ingestion Layer \timing{5:20 -- 5:45}}

\begin{scriptbox}
``Tóm lại phần Ingestion: hệ thống có \textbf{hai luồng vào song song}:
luồng batch qua CSV -- PyIceberg, và luồng streaming qua Redpanda -- Consumer.
Cả hai được ghi vào cùng Bronze Iceberg table, không conflict nhau,
vì Iceberg đảm bảo ACID.

Prefect điều phối toàn bộ luồng với quality gate Soda Core ngăn
data xấu tiếp tục lên Silver. Mỗi task có retry tự động.

Đáng chú ý là toàn bộ phần này \textbf{không có file YAML cấu
hình} --- mọi thứ là Python thuần, rất dễ test và debug.''
\end{scriptbox}

\subsection*{Chuyển tiếp sang Người 3 \timing{5:45 -- 6:00}}

\begin{scriptbox}
``Dữ liệu đã vào Bronze. Bước tiếp theo là \textbf{transform} từ
Bronze lên Silver rồi Gold --- và \textbf{lineage} để biết dữ liệu
đã đi qua những bước nào. Mình xin nhường lại lời cho bạn [tên Người 3].''
\end{scriptbox}

\end{speakerbox}

\newpage

% ============================================================
\section{Câu hỏi thường gặp --- Q\&A Chuẩn bị}
% ============================================================

\begin{tcolorbox}[colback=gray!5, colframe=gray!50, arc=4pt,
  title={\faQuestion\ 7 câu hỏi thầy cô hay hỏi}]

\textbf{Q1: Tại sao dùng Redpanda mà không dùng Kafka thật?}

\begin{scriptbox}
``Redpanda Kafka-compatible hoàn toàn --- code Python không cần thay đổi.
Điểm khác biệt: Redpanda chạy trong 1 container, không cần Zookeeper,
không cần JVM --- tiêu thụ dưới 500~MB RAM so với Kafka cần 2~GB+.
Với quy mô đồ án sinh viên, Redpanda giúp deploy đơn giản hơn rất nhiều.''
\end{scriptbox}

\textbf{Q2: Tại sao Bronze append-only? Nếu source data sai thì sao?}

\begin{scriptbox}
``Bronze là \textbf{immutable raw data} --- giữ lại đúng những gì hệ thống
ngoài gửi vào. Nếu source sai, chúng ta append bản sửa lên Bronze,
rồi Silver MERGE INTO sẽ dùng bản mới nhất. Bronze giữ toàn bộ
lịch sử để audit.''
\end{scriptbox}

\textbf{Q3: Nếu Consumer chết giữa chừng thì message có mất không?}

\begin{scriptbox}
``Không. Consumer chỉ \code{commit offset} sau khi đã \code{flush} thành
công vào Iceberg. Nếu crash trước commit, lần restart tiếp theo
consumer sẽ đọc lại từ offset cũ --- \textbf{at-least-once delivery}.
Iceberg ACID loại duplicate nếu cần.''
\end{scriptbox}

\textbf{Q4: Prefect khác gì Airflow?}

\begin{scriptbox}
``Prefect dùng Python decorator \code{@flow} / \code{@task} --- viết như
code Python bình thường, không cần DAG class hay YAML cấu hình.
Prefect Cloud có free tier, mới nên không cần tự host UI.
Airflow nặng hơn, cần server riêng, cú pháp DAG phức tạp hơn.''
\end{scriptbox}

\textbf{Q5: Quality gate Soda Core kiểm tra cụ thể những gì?}

\begin{scriptbox}
``Ba loại check chính:
(1) \textbf{Completeness}: not\_null trên các khóa chính (\code{order\_id}, \code{customer\_id}).
(2) \textbf{Volume}: row\_count phải \code{>= 90000} --- phát hiện trường hợp mất data.
(3) \textbf{Referential}: mỗi \code{order\_id} trong \code{order\_items} phải có trong \code{orders}.''
\end{scriptbox}

\textbf{Q6: Tại sao không dùng dbt thay vì PySpark cho Silver/Gold?}

\begin{scriptbox}
``dbt đã bị loại khỏi dự án. Lý do: dbt cần adapter riêng cho
Iceberg, thêm toolchain phức tạp. PySpark đọc/ghi Iceberg native qua
\code{spark-iceberg} extension, cùng engine xử lý cả small và TB-scale.
Không cần thêm dependency.''
\end{scriptbox}

\textbf{Q7: Hệ thống này scale lên production thật được không?}

\begin{scriptbox}
``Được. Thay Cloudflare R2 bằng AWS S3, thêm Spark cluster (master +
workers) thay cho \code{local[2]}, thêm Kafka cluster thay Redpanda. Các
abstraction lớp trên (PyIceberg, confluent-kafka, PySpark) giữ nguyên ---
không thay đổi code.''
\end{scriptbox}

\end{tcolorbox}

\newpage

% ============================================================
\section{Cheat Sheet --- Tóm tắt 1 trang}
% ============================================================

\begin{multicols}{2}

\begin{keybox}{\faBolt\ Batch Ingestion}
\begin{itemize}[leftmargin=1em, itemsep=2pt, topsep=2pt]
  \item 8 CSV $\to$ Pandas $\to$ PyIceberg
  \item \code{table.append(pa\_table)} --- ACID
  \item \code{@task(retries=3)} --- tự retry
  \item Bronze: append-only, immutable
  \item 10 Iceberg tables trên R2
\end{itemize}
\end{keybox}

\columnbreak

\begin{keybox}{\faRss\ Redpanda Streaming}
\begin{itemize}[leftmargin=1em, itemsep=2pt, topsep=2pt]
  \item 1 container, không Zookeeper
  \item Kafka-compatible API
  \item \textless500~MB RAM vs Kafka 2~GB
  \item Speed factor 86400 (demo)
  \item Key = order\_id (partition order)
  \item Micro-batch flush 50 msgs
\end{itemize}
\end{keybox}

\end{multicols}

\begin{multicols}{2}

\begin{keybox}{\faTasks\ Prefect Orchestration}
\begin{itemize}[leftmargin=1em, itemsep=2pt, topsep=2pt]
  \item \code{@flow} + \code{@task} --- Python thuần
  \item \code{wait\_for} = dependency graph
  \item \code{retries=3} mỗi task
  \item Schedule Daily 02:00 UTC
  \item Prefect Cloud free tier UI
\end{itemize}
\end{keybox}

\columnbreak

\begin{keybox}{\faShieldAlt\ Soda Core Quality Gate}
\begin{itemize}[leftmargin=1em, itemsep=2pt, topsep=2pt]
  \item Check sau mỗi layer Bronze/Silver
  \item Null check, row count, referential
  \item Fail $\to$ PAUSE flow, không lên Silver
  \item Webhook notification khi fail
\end{itemize}
\end{keybox}

\end{multicols}

\vspace{0.8em}

% ── Mini pipeline diagram ─────────────────────────────────────
\begin{center}
\begin{tikzpicture}[
  node distance=0.2cm and 0.5cm,
  s/.style={draw, rounded corners=3pt, minimum width=1.9cm,
    minimum height=0.55cm, align=center, font=\small, line width=0.8pt},
  a/.style={-Stealth, line width=0.8pt},
  g/.style={diamond, draw=orange!60, fill=yellow!20, aspect=1.8,
    minimum height=0.45cm, font=\tiny, align=center},
]
\node[s, fill=gray!15, draw=gray]       (csv)   {CSV};
\node[s, fill=orange!20, draw=orange!70, right=0.5cm of csv]  (rp)  {Redpanda};
\node[s, fill=orange!25, draw=orange!80, right=0.5cm of rp]   (bz)  {Bronze\\Iceberg};
\node[g, right=0.5cm of bz]                                   (qb)  {Soda\\OK?};
\node[s, fill=silver!40, draw=silver!80, right=0.5cm of qb]   (sv)  {Silver\\PySpark};
\node[g, right=0.5cm of sv]                                   (qs)  {Soda\\OK?};
\node[s, fill=gold!40, draw=gold!90,    right=0.5cm of qs]    (gd)  {Gold\\PySpark};
\node[s, fill=green!15, draw=green!60,  right=0.5cm of gd]    (out) {BI / ML};

\draw[a] (csv) -- (rp);
\draw[a] (rp)  -- (bz);
\draw[a] (bz)  -- (qb);
\draw[a] (qb)  -- node[font=\tiny,above]{Pass} (sv);
\draw[a] (sv)  -- (qs);
\draw[a] (qs)  -- node[font=\tiny,above]{Pass} (gd);
\draw[a] (gd)  -- (out);

\node[font=\tiny, color=gray,      below=0.08cm of bz] {append-only};
\node[font=\tiny, color=gray,      below=0.08cm of sv] {MERGE INTO};
\node[font=\tiny, color=gray,      below=0.08cm of gd] {partitioned};
\end{tikzpicture}
\end{center}

\vspace{1em}
\begin{center}
\small\color{gray}
Script Người 2 --- Olist Data Lakehouse \textbullet\
Liên quan: \code{streaming/producer.py} \textbullet\
\code{streaming/consumer.py} \textbullet\
\code{prefect/flows/full\_pipeline.py} \textbullet\
\code{quality/soda\_runner.py}
\end{center}

% ============================================================
\end{document}
% ============================================================
